% SPDX-FileCopyrightText: 2026 Yiyu Zhou <yiyu@yiyuzhou.io>
%
% SPDX-License-Identifier: GPL-3.0-or-later

\chapter{Sequences}
\label{ch:sequences}

So far, each cue has needed its own GO\@.  A cue can also fire the next cue
by itself, so that a single GO plays a whole sequence.  Three settings on the
\ui{Basics} tab control the timing:

\begin{description}
\item[Pre-wait] Seconds between the cue firing and its action starting.
\item[Post-wait] Seconds between the action starting and the next cue
  firing, for a cue that continues to the next.
\item[Then] What happens to the next cue: \ui{wait for GO}, the default,
  leaves it to the operator; \ui{auto-continue} fires it when the post-wait
  ends; and \ui{auto-follow} fires it when this cue's action ends.
\end{description}

The \texttt{Pre}, \texttt{Post}, and \texttt{Then} columns of the cue list
show the same settings.

\section{A sequence in the storm}

Press \key{Control+l} and open \texttt{storm.cinq.toml} again.

\begin{enumerate}
\item Select cue~1, press \ui{Standby selected}, and press \ui{GO} to restart
  the preshow music.  Cue~3 is now on standby, and the transport summarizes
  what its GO fires: three cues, the second 4.0 seconds and the third 5.5
  seconds after GO, as Figure~\ref{fig:playing} in
  Chapter~\ref{ch:running} shows.
\item Press \ui{GO}\@.  Cue~3 fades the music out over six seconds.  Four
  seconds after GO, the distant thunder starts, followed 1.5 seconds later by
  the rain, as in Figure~\ref{fig:running}.
\end{enumerate}

\begin{figure}[H]
  \centering
  \includegraphics[width=\textwidth]{screenshots/running.png}
  \caption{Five seconds after GO on cue 3, which fades out the preshow
    music, while the thunder plays and the rain is about to start}
  \label{fig:running}
\end{figure}

Cue~3 auto-continues to cue~4 after a post-wait of 4 seconds, and cue~4
continues to cue~5 after 1.5 seconds.  Cue~5 waits for GO, so the sequence
ends there.  While a cue waits to start, the sidebar counts down to it.

\section{Pre-wait and post-wait}

A pre-wait can produce the same timing, but it belongs to the cue that waits
rather than to the cue before it.

\begin{enumerate}
\item Press \ui{Panic}.  Set the \ui{Post-wait} of cue~3 to 0 and the
  \ui{Pre-wait} of cue~4 to 4.
\item Put cue~1 on standby, press \ui{GO} for the music, and press \ui{GO}
  again for cue~3.  The sequence sounds the same: cue~4 now fires at once and
  waits four seconds before playing.
\item Put cue~4 on standby and press \ui{GO}\@.  The thunder now waits four
  seconds; before the change, it would have started immediately.  A pre-wait
  delays its own cue's action however the cue is fired, while a post-wait
  delays only the next cue.
\item Restore the original values: a post-wait of 4 on cue~3 and a pre-wait
  of 0 on cue~4.
\end{enumerate}

\section{Auto-follow}

Auto-follow fires the next cue when the current cue's action ends, such as
when a sound finishes playing, so it needs no post-wait.  An action that
never ends, such as a looping sound or a picture, never triggers an
auto-follow.

\section{Wait and memo cues}

A wait cue lasts for its \ui{Duration} and does nothing else; it paces a
sequence.  A memo cue does nothing at all.  It serves as a note in the list
or as the start of a sequence, like cue~1 of the storm show, whose notes tell
the operator what the music is for.

\section{Pause and start cues}

A pause cue pauses another cue, its target, and a start cue starts it again.
Both take no time, and both go in the list like any other cue.

Press \key{Control+l} and open \texttt{my-show.cinq.toml}.

\begin{enumerate}
\item Select the rain.  On its \ui{Time \& Loops} tab, check \ui{Loop until
  stopped or devamped} if it is not checked, so that the rain keeps playing
  while you work.
\item Choose \ui{Pause} in the \ui{New cue} list of the toolbar and press
  \ui{Add cue}\@.  The new cue appears below the rain and targets it, since a
  new pause cue targets the cue above it.
\item With the pause cue selected, choose \ui{Start} in the \ui{New cue} list
  and press \ui{Add cue}\@.  The new start cue appears below the pause cue,
  which it targets for the same reason.  Open the \ui{Target} list on its
  \ui{Basics} tab and choose the rain instead.
\item Select the rain, press \ui{Standby selected}, and press \ui{GO}\@.  The
  rain starts.  Press \ui{GO} again.  The pause cue pauses the rain, which
  the \texttt{State} column and the sidebar mark as paused, and the sound
  stops where it was.
\item Press \ui{GO} again.  The start cue resumes the rain from the point
  where it stopped.  \ui{Resume all} in the transport resumes it as well.
\end{enumerate}

A start cue whose target is not paused fires the target as GO would, together
with the cues that chain to it.  A pause cue also holds the cues that wait
for its target to start them, as cue~4 of the storm show waits out the
post-wait of cue~3, so that the whole sequence holds where it is.  Neither
moves the cue on standby.

A GoTo cue does move it.  Its target can be any cue, and when GO fires the
GoTo cue, the target goes on standby, so that the next GO fires it, and
nothing else happens.  Use one at the end of a list to put its first cue on
standby again.

\section{Loop sections and devamp cues}

A loop need not take the whole file.  \ui{Loop start} and \ui{Loop end}, on
the \ui{Time \& Loops} tab beside \ui{Start} and \ui{End}, set the part of the
file that a looping cue plays again and again, in seconds from the beginning
of the file.  The cue plays from its start to its loop end, then goes back to
its loop start each time it reaches the loop end.  Left blank, they stand for
the start and the end of the cue, so a cue without them loops its whole file,
as before.  The fields are available only while the cue loops.  One file can
therefore hold an introduction, a vamp that repeats under a scene, and an
ending.

A devamp cue lets the vamp end.  Its target, an audio or video cue that
loops, plays on to the end of its current pass through the loop, goes on from
the loop end to the end of the file, and ends there.  Like a pause cue, a
devamp cue takes no time.

Press \key{Control+l} and open \texttt{my-show.cinq.toml}.  The rain plays for
eight seconds, so the steps make its first two seconds the introduction, the
next four the vamp, and the last two the ending.

\begin{enumerate}
\item Press \ui{Panic}, and select the rain.  On its \ui{Time \& Loops} tab,
  check \ui{Loop until stopped or devamped} if it is not checked, and type
  \texttt{2} in \ui{Loop start} and \texttt{6} in \ui{Loop end}.
\item With the rain selected, choose \ui{Devamp} in the \ui{New cue} list
  and press \ui{Add cue}\@.  The new cue appears below the rain and targets
  it, as a new pause cue does.
\item Select the rain, press \ui{Standby selected}, and press \ui{GO}\@.  The
  rain plays its two-second introduction, then repeats seconds 2 to 6 of the
  file, and the \texttt{State} column reads \texttt{looping}.  The devamp cue
  is on standby.
\item Press \ui{GO} while the rain loops.  The \texttt{State} column of the
  rain shows the time it has left instead of \texttt{looping}.  The rain
  finishes its pass, plays the last two seconds of its file, and ends.
\end{enumerate}

A devamp cue never cuts a pass short.  It lets the rain leave at the first
end of a pass that comes more than half a second after the cue acts, so that a
GO that comes just before an end waits for the end after it.  A target that
does not loop cannot be devamped: its devamp cue reads \texttt{NO LOOP}, and GO
skips it.  A loop that starts at or after its end cannot play either, and its
cue reads \texttt{BROKEN}\@.

\section*{Exercise}

You can now make one GO fire a whole sequence.  In your own show, move the
thunder cue from Chapter~\ref{ch:cues} above the rain with
\key{Alt+ArrowUp}, and set its \ui{Then} to \ui{auto-follow}, so that one GO
plays the thunder and then the rain.
